% GENERATED FILE. Edit canonical lessons, then run npm run generate:books.
% canonical-source-hash: fnv1a64:cf0d2223aff35ea2
% canonical-lessons: MW-C77-lal, MW-C77-nilo, MW-C77-kalo, MW-C77-bhuro, MW-C77-mitho

\chapter{Red, Blue, Black, Brown, Sweet}
\label{ch:mw-red-blue-black-brown-sweet}
\hlchaptermodality{77}

\begin{chapteropening}
\textbf{By the end of this chapter:} \emph{I can say red, blue, black, brown and sweet.}

Complete the last lesson of chapter 77: I can say red, blue, black, brown and sweet.
\end{chapteropening}

\section[\texorpdfstring{\mw{लाल} (lāl)}{लाल (lāl)}]{\mw{लाल} (lāl) — red}

\label{lesson:MW-C77-lal}



\begin{quote}
Before the new one: say the Marwadi for green, then the Marwadi for yellow.
\end{quote}

\subsection*{You'll want to know: \mw{लाल}}
\textbf{\mw{लाल}} — \emph{lāl} — \textquotedblleft{}red\textquotedblright{}.

\begin{cousinweb}
\textbf{\mw{लाल}} (\emph{lāl}) is a loan, usually traced to Persian \emph{lāl}, \textquotedblleft{}red; ruby\textquotedblright{}.
\end{cousinweb}

\begin{grammarlens}[title={What you've built}]
One more word for this chapter.
\end{grammarlens}

\subsection*{Your turn}
\begin{itemize}
\raggedright
  \item \emph{Say it:} \emph{lāl}
  \item \emph{Say it:} \emph{lāl}, once more
  \item \emph{Say it:} say \emph{lāl}
  \item \emph{Recall:} say the Marwadi for green, then the Marwadi for yellow, then say \emph{lāl} again
\end{itemize}

\begin{tcolorbox}[breakable,colback=teal!4,colframe=teal!35!black,title={Before you move on}]
What does \mw{लाल} mean? (\textquotedblleft{}Red\textquotedblright{}.) Say it once more.
\end{tcolorbox}
\section[\texorpdfstring{\mw{नीलो} (nīlo)}{नीलो (nīlo)}]{\mw{नीलो} (nīlo) — blue}

\label{lesson:MW-C77-nilo}



\begin{quote}
Before the new one: say the Marwadi for yellow, then the Marwadi for red.
\end{quote}

\subsection*{You'll want to know: \mw{नीलो}}
\textbf{\mw{नीलो}} — \emph{nīlo} — \textquotedblleft{}blue\textquotedblright{}.

\begin{cousinweb}
\textbf{\mw{नीलो}} (\emph{nīlo}) is a tadbhava, a word inherited from Sanskrit and reshaped by sound change. It comes from \emph{nīla}, \textquotedblleft{}dark blue, dark green, black\textquotedblright{}. CDIAL gives no Marwari form in this entry (CDIAL 7563); the match is through Hindi \emph{nīlā} and Gujarati \emph{nīlũ}.
\end{cousinweb}

\begin{grammarlens}[title={What you've built}]
One more word for this chapter.
\end{grammarlens}

\subsection*{Your turn}
\begin{itemize}
\raggedright
  \item \emph{Say it:} \emph{nīlo}
  \item \emph{Say it:} \emph{nīlo}, once more
  \item \emph{Say it:} say \emph{nīlo}
  \item \emph{Recall:} say the Marwadi for yellow, then the Marwadi for red, then say \emph{nīlo} again
\end{itemize}

\begin{tcolorbox}[breakable,colback=teal!4,colframe=teal!35!black,title={Before you move on}]
What does \mw{नीलो} mean? (\textquotedblleft{}Blue\textquotedblright{}.) Say it once more.
\end{tcolorbox}
\section[\texorpdfstring{\mw{काळो} (kāḷo)}{काळो (kāḷo)}]{\mw{काळो} (kāḷo) — black}

\label{lesson:MW-C77-kalo}



\begin{quote}
Before the new one: say the Marwadi for red, then the Marwadi for blue.
\end{quote}

\subsection*{You'll want to know: \mw{काळो}}
\textbf{\mw{काळो}} — \emph{kāḷo} — \textquotedblleft{}black\textquotedblright{}.

\begin{cousinweb}
\textbf{\mw{काळो}} (\emph{kāḷo}) is a tadbhava, a word inherited from Sanskrit and reshaped by sound change. It comes from \emph{kāla}, \textquotedblleft{}black, dark-blue\textquotedblright{}. CDIAL lists Marwari \emph{kāḷo} under it (CDIAL 3083). CDIAL reports a view that the Sanskrit word came from Dravidian.
\end{cousinweb}

\begin{grammarlens}[title={What you've built}]
One more word for this chapter.
\end{grammarlens}

\subsection*{Your turn}
\begin{itemize}
\raggedright
  \item \emph{Say it:} \emph{kāḷo}
  \item \emph{Say it:} \emph{kāḷo}, once more
  \item \emph{Say it:} say \emph{kāḷo}
  \item \emph{Recall:} say the Marwadi for red, then the Marwadi for blue, then say \emph{kāḷo} again
\end{itemize}

\begin{tcolorbox}[breakable,colback=teal!4,colframe=teal!35!black,title={Before you move on}]
What does \mw{काळो} mean? (\textquotedblleft{}Black\textquotedblright{}.) Say it once more.
\end{tcolorbox}
\section[\texorpdfstring{\mw{भूरो} (bhūro)}{भूरो (bhūro)}]{\mw{भूरो} (bhūro) — brown}

\label{lesson:MW-C77-bhuro}



\begin{quote}
Before the new one: say the Marwadi for blue, then the Marwadi for black.
\end{quote}

\subsection*{You'll want to know: \mw{भूरो}}
\textbf{\mw{भूरो}} — \emph{bhūro} — \textquotedblleft{}brown\textquotedblright{}.

\begin{cousinweb}
\textbf{\mw{भूरो}} (\emph{bhūro}) is a tadbhava, a word inherited and reshaped by sound change. It comes from \emph{*bhrūra}, \textquotedblleft{}brown\textquotedblright{}, a form CDIAL reconstructs from its descendants. CDIAL gives no Marwari form in this entry (CDIAL 9690); the match is through Hindi \emph{bhūrā} and Gujarati \emph{bhūrũ}. CDIAL sets it beside Indo-European \emph{*bhrūno-}, behind English \emph{brown}.
\end{cousinweb}

\begin{grammarlens}[title={What you've built}]
One more word for this chapter.
\end{grammarlens}

\subsection*{Your turn}
\begin{itemize}
\raggedright
  \item \emph{Say it:} \emph{bhūro}
  \item \emph{Say it:} \emph{bhūro}, once more
  \item \emph{Say it:} say \emph{bhūro}
  \item \emph{Recall:} say the Marwadi for blue, then the Marwadi for black, then say \emph{bhūro} again
\end{itemize}

\begin{tcolorbox}[breakable,colback=teal!4,colframe=teal!35!black,title={Before you move on}]
What does \mw{भूरो} mean? (\textquotedblleft{}Brown\textquotedblright{}.) Say it once more.
\end{tcolorbox}
\section[\texorpdfstring{\mw{मीठो} (mīṭho)}{मीठो (mīṭho)}]{\mw{मीठो} (mīṭho) — sweet}

\label{lesson:MW-C77-mitho}



\begin{quote}
Before the new one: say the Marwadi for black, then the Marwadi for brown.
\end{quote}

\subsection*{You'll want to know: \mw{मीठो}}
\textbf{\mw{मीठो}} — \emph{mīṭho} — \textquotedblleft{}sweet\textquotedblright{}.

\begin{cousinweb}
\textbf{\mw{मीठो}} (\emph{mīṭho}) is a tadbhava, a word inherited from Sanskrit and reshaped by sound change. It comes from \emph{miṣṭa}, \textquotedblleft{}dainty, sweet\textquotedblright{}. CDIAL lists Old Marwari \emph{sumīṭha} under it (CDIAL 10299). The Old Marwari row means \textquotedblleft{}very sweet\textquotedblright{}; Hindi has \emph{mīṭhā}.
\end{cousinweb}

\begin{grammarlens}[title={What you've built}]
One more word for this chapter.
\end{grammarlens}

\subsection*{Your turn}
\begin{itemize}
\raggedright
  \item \emph{Say it:} \emph{mīṭho}
  \item \emph{Say it:} \emph{mīṭho}, once more
  \item \emph{Say it:} say \emph{mīṭho}
  \item \emph{Recall:} say the Marwadi for black, then the Marwadi for brown, then say \emph{mīṭho} again
\end{itemize}

\begin{tcolorbox}[breakable,colback=teal!4,colframe=teal!35!black,title={Before you move on}]
What does \mw{मीठो} mean? (\textquotedblleft{}Sweet\textquotedblright{}.) Say it once more.
\end{tcolorbox}
